\section{The Construction}

Let $N^{2n}$ be any $2n$-dimensional totally non-spin manifold, $n\geq 2$. We could for example take $N:= \CP^2 \times T^{2n-4}$. Consider the manifold
\[
M := \CP^n \# T^{2n} \# N.
\]
Since the collapse map $M \to T^{2n}$ has degree $1$ and $T^{2n}$ is enlargeable, it follows from Lemma \ref{lem:constructing_enlargeable_mnf} that $M$ is also enlargeable. We will now construct a circle bundle $E \to M$ as follows. Let 
\[
c_1 \in H^2(M;\Z) \cong H^2(\CP^n;\Z) \oplus H^2(T^{2n};\Z) \oplus H^2(N;\Z)
\]
be the cohomology class corresponding to the first chern class $\alpha \in H^2(\CP^n;\Z)$ of the canonical bundle $S^{2n+1} \to \CP^n$. We then define $E \to M$ as the unique $S^1$-principal bundle with first chern class equal to $c_1$. Now this is quite an abstract definition, so let us develop an intuition for how this bundle looks like. We claim that $E$ is trivial over the summand $T^{2n} \# N$ and is given by the restriction of the canonical bundle $S^{2n+1} \to \CP^n$ over the $\CP^n$ summand. To make this precise let 
\[
i_1 \colon (T^{2n} \# N) - D^{2n} \to M
\]
be the inclusion. Then
\[\begin{tikzcd}
	{H^2(M; \Z)} & {H^2((T^{2n} \# N) - D^{2n};\Z) } \\
	{H^2(\CP^n;\Z) \oplus H^2(T^{2n} \# N;\Z)} & {H^2(T^{2n} \# N;\Z) }
	\arrow["{i_1^\ast}", from=1-1, to=1-2]
	\arrow["\cong", from=2-1, to=1-1]
	\arrow["{\mathrm{proj.}}", from=2-1, to=2-2]
	\arrow["\cong", from=2-2, to=1-2]
\end{tikzcd}\]
commutes and therefore $i_1^\ast (c_1) = 0$. Again using the fact that $S^1$-principal bundles are uniquely determined by their first chern class we see that $i_1^\ast E$ is indeed trivial. Similiarly if we denote by 
\[
i_2 \colon \CP^n - D^{2n} \to M, \quad j \colon \CP^{2n}-D^{2n} \to \CP^{2n}
\]
the inclusions, the diagram
\[\begin{tikzcd}
	{H^2(M ;\Z)} & {H^2(\CP^n - D^{2n};\Z) } \\
	{H^2(\CP^n;\Z) \oplus H^2(T^{2n} \# N;\Z)} & {H^2(\CP^n;\Z) }
	\arrow["{i_2^\ast}", from=1-1, to=1-2]
	\arrow["\cong", from=2-1, to=1-1]
	\arrow["{\mathrm{proj.}}", from=2-1, to=2-2]
	\arrow["\cong", "j^\ast"', from=2-2, to=1-2]
\end{tikzcd}\]
commutes and thus $i_2^\ast(c_1) = j^\ast(\alpha)$. So the first chern classes of $i_2^\ast E$ and $j^\ast S^{2n+1}$ agree, hence they must be isomorphic bundles. 

With this intuition in mind we now want to show the existence of a metric of positive scalar curvature on $E$ by applying Theorem \ref{thm:gromov_hanke}. In order to apply the Theorem we will show the following two claims.

\begin{claim}
	\label{claim1}
$E$ is totally non-spin.
\end{claim}

\begin{claim}
	\label{claim2}
The projection $p \colon E \to M$ induces an isomorphism on $\pi_1$.
\end{claim}

Granting these Claims for now let us see why they suffice for the assumptions of Theorem \ref{thm:gromov_hanke} to be met.  
For this we would like to show that $E$ is null-homologous as a cycle in the group homology $H_{2n}(\pi_1(E);\Z)$, which can be seen by using Claim \ref{claim2} as follows. Let $q \colon DE \to M$ be the disk bundle corresponding to $E$. Moreover let $\phi \colon E \to B\pi_1(E)$ and $\Phi \colon DE \to B\pi_1(DE)$ be the classifiying maps of the respective universal coverings. Then the diagram 
\[\begin{tikzcd}
	DE & {B\pi_1(DE)} \\
	E & {B\pi_1(E)}
	\arrow["\Phi", from=1-1, to=1-2]
	\arrow["\iota", from=2-1, to=1-1]
	\arrow["\phi", from=2-1, to=2-2]
	\arrow["{B\pi_1(\iota)}"', from=2-2, to=1-2]
\end{tikzcd}\]
commutes (Begründung rein!), where the vertical maps are (induced by) the inclusion $\iota \colon E \to DE$. We claim that $B\pi_1(\iota)$ is an isomorphism. This follows from commutativity of  
\[\begin{tikzcd}[column sep = small]
	{B\pi_1(E)} && {B\pi_1(DE)} \\
	& {B\pi_1(M)}
	\arrow["{B\pi_1(\iota)}", from=1-1, to=1-3]
	\arrow["{B\pi_1(p)}"',"\cong", from=1-1, to=2-2]
	\arrow["{B\pi_1(q)}","\cong"', from=1-3, to=2-2]
\end{tikzcd}\]
where $B\pi_1(q)$ is an isomorphism by homotopy invariance and $B\pi_1(p)$ is an isomorphism by Claim \ref{claim2}. Thus $\phi$ factors over $\iota$. Since $DE$ is an oriented null-cobordism of $E$ it follows that $\iota_\ast[E] = 0 \in H_{2n}(DE;\Z)$ and so in particular $\phi_\ast [E] = 0\in H_{2n}(\pi_1(E);\Z)$ as we wanted to show. So together with Claim \ref{claim1} all assumptions of Theorem \ref{thm:gromov_hanke} are met, showing that $E$ admits a metric of positive scalar curvature. We will now proof Claim \ref{claim1} and Claim \ref{claim2}

\begin{proof}[Proof of Claim \ref{claim1}]
	Let $i \colon N- D^{2n} \to M$ be the inclusion. We have already established above, that $E$ is trivial over the $T^{2n} \# N$ summand and so in particular
	\[
	i^\ast E \cong (N- D^{2n}) \times S^1.
	\]
	Consequently the universal cover of $i^\ast E$ is given by $\tilde{N} \times \R$, where $\tilde{N}$ is the universal cover of $N$. Let $\pi_{\tilde{N}} \colon \tilde{N} \times \R \to \tilde{N}$ and $\pi_\R \colon \tilde{N} \times \R \to \R$ be the projections. We compute 
	\[
	w_2(T(\tilde{N} \times \R)) = w_2(\pi_{\tilde{N}}^\ast T\tilde{N} \oplus \pi_\R^\ast T\R) = \pi_{\tilde{N}}^\ast w_2(T\tilde{N}).
	\]
	Since $\pi_{\tilde{N}}$ is a homotopy equivalence and $\tilde{N}$ is by construction non-spin, it follows that $w_2(T(\tilde{N} \times \R)) \neq 0$, such that $i^\ast E$ is totally non-spin. Therefore $M$ contains a totally non-spin submanifold of codimension $0$ and is consequently itself totally non-spin by (Lemma in Preliminaries section einfügen)
\end{proof}

\begin{proof}[Proof of Claim \ref{claim2}]
	Let $j \colon \CP^{n} - D^{2n}$ be the inclusion. Consider the following commutative diagram, where the rows are taken from the long exact sequences of homotopy groups of the respective fibrations and the vertical maps are induced by inclusions: 
	\[\begin{tikzcd}
	{\pi_2(M)} & {\pi_1(S^1)} & {\pi_1(E)} \\
	{\pi_2(\CP^n - D^{2n})} & {\pi_1(S^1)} & {\pi_1(j^\ast E)} \\
	{\pi_2(\CP^n)} & {\pi_1(S^1)} & {\pi_1(S^{2n+1})}
	\arrow["{\partial_1}", from=1-1, to=1-2]
	\arrow[from=1-2, to=1-3]
	\arrow[from=2-1, to=1-1]
	\arrow["{\partial_2}", from=2-1, to=2-2]
	\arrow["\cong"', from=2-1, to=3-1]
	\arrow["{\mathrm{id}}", from=2-2, to=1-2]
	\arrow[from=2-2, to=2-3]
	\arrow["{\mathrm{id}}"', from=2-2, to=3-2]
	\arrow[from=2-3, to=1-3]
	\arrow[from=2-3, to=3-3]
	\arrow["{\partial_3}", from=3-1, to=3-2]
	\arrow[from=3-2, to=3-3]
\end{tikzcd}\]
The vertical map on the lower left is an isomorphism, because both $\CP^n$ and $\CP^n - D^{2n}$ are simply connected and so by the Hurewicz Theorem
\[\begin{tikzcd}
	{\pi_2(\CP^n-D^{2n})} & {\pi_2(\CP^n)} \\
	{H_2(\CP^n-D^{2n};\Z)} & {H_2(\CP^n;\Z)}
	\arrow[from=1-1, to=1-2]
	\arrow["\cong", from=1-1, to=2-1]
	\arrow["\cong", from=1-2, to=2-2]
	\arrow["\cong", from=2-1, to=2-2]
\end{tikzcd}\]
commutes and the vertical maps are isomorphisms. We now claim that $\partial_1$ is surjective. For this it suffices to show surjectivity of $\partial_2$, which in turn is equivalent to surjectivity of $\partial_3$. But this is clear by exactness of the rows and $\pi_1(S^{2n+1}) = 0$. Consequently
\[
\pi_1(S^1) \to \pi_1(E)
\]
is the zero map and therefore we conclude from exactness of
\[
\pi_1(S^1) \to \pi_1(E) \to \pi_1(M) \to \pi_0(S^1),
\]
that $\pi_1(E) \to \pi_1(M)$ is indeed an isomorphism.
\end{proof}

The above construction gives examples to Question \ref{the_question} in all even dimensions $\geq 4$. To also get examples in odd dimensions let $E \to M$ be as above and let $\pi \colon M \times S^1 \to M$ be the projection. Then the pullback bundle $\pi^\ast E \cong E \times S^1$ clearly admits a metric of positive scalar curvature and by Lemma \ref{lem:constructing_enlargeable_mnf} the base space $M \times S^1$ is still enlargeable.